\subsection*{Tyrosine catabolism is enriched under CTL, and the dhurrin route is recovered under LIN}

Candidates from both trials cluster around a single metabolic branch point on opposite sides. In sorghum, tyrosine is either catabolized via tyrosine aminotransferase (TAT) and 4-hydroxyphenylpyruvate dioxygenase (HPPD) to homogentisate and then to fumarate and acetoacetate, supplying acetyl-CoA for fatty-acid synthesis \citep{Schenck2018Tyrosine}, or diverted to dhurrin biosynthesis via CYP79A1, CYP71E1, and UGT85B1 \citep{Busk2002Dhurrin,Darbani2016Dhurrin}. Because both pathways use the same precursor, candidates on either side suggest the trials differ in tyrosine allocation.

Under CTL, the \textit{L-tyrosine catabolic process} is enriched, and all associated lipids are TGs. Two of the four genes carrying the term are consecutive enzymes of the pathway, the other two being nicotianamine aminotransferases \citep{Inoue2008NAAT}. \texttt{SORBI\_3005G152200} is the only tyrosine aminotransferase candidate and is associated with ten TGs, eight of which carry palmitate (16:0), and the next enzyme, 4-hydroxyphenylpyruvate dioxygenase (\texttt{SORBI\_3002G104200}), is a candidate for a single lipid, TG(10:0\_10:0\_10:0). They sit on different chromosomes, so the term draws on independent loci rather than on one LD block. The pathway ends in fumarate and acetoacetate \citep{Schenck2018Tyrosine}, and acetoacetate is converted to acetyl-CoA, the carbon substrate for \textit{de novo} fatty acid synthesis, with those fatty acids then esterified onto the glycerol backbone to form TG \citep{Yen2008DGAT}. Tyrosine catabolism therefore offers a plausible, if indirect, route from this locus to genotypic variation in TG content.

Under LIN, no tyrosine- or dhurrin-related GO term is enriched. Dhurrin is instead implicated by a single chromosome 1 locus that is strongly associated. The seven candidates share 12/13 PA-to-class ratios; three also carry the fatty-acid class sum and FA(22:1), and four carry TG(14:0\_18:3\_18:3). Three candidates are known dhurrin-cluster genes \citep{Darbani2016Dhurrin}—UGT85B1 (\texttt{SORBI\_3001G012400}), SbMATE2 (\texttt{SORBI\_3001G012600}), and SbGST1 (\texttt{SORBI\_3001G012500})—and a fourth, \texttt{SORBI\_3001G012000}, is a CYP71 P450 annotated to the oxime monooxygenase step. The association does not extend to CYP79A1 (\texttt{SORBI\_3001G012300}) or CYP71E1 (\texttt{SORBI\_3001G012200}); although they lie between candidates, they are not tagged at $r^2\ge 0.4$, indicating the signal covers only part of the cluster. Because all seven candidates share the same signal, the association maps a region rather than a specific gene. Only two candidates appear in any (broad) enriched term—iron ion binding (P450) and transmembrane transporter activity (transporter)—while UGT85B1 and SbGST1 appear in none. Thus, the dhurrin link is based on locus gene content and annotation, not enrichment.

Dhurrin is an alternative fate of tyrosine and functions as both a nitrogen store and a defense compound, making it relevant under reduced nitrogen. Nitrogen supply induces dhurrin synthesis in sorghum \citep{Busk2002Dhurrin}, and its turnover releases nitrogen for early growth. High-dhurrin genotypes accumulate more shoot biomass than low-dhurrin genotypes under nitrogen deficiency \citep{Emendack2025Dhurrin}. Abiotic stress can also increase dhurrin independently of nitrogen. Salt stress upregulates tyrosine-metabolism genes and coincides with growth restriction and dhurrin accumulation \citep{Amombo2026Tyrosine}. Combined water, salinity, and heat stress doubles dhurrin while shifting the forage lipid profile \citep{GiraldoAcosta2026Dhurrin}, and dhurrin contributes to pathogen defense \citep{Fall2025Dhurrin}. Under low inputs, reduced commitment of tyrosine-derived carbon and nitrogen to dhurrin, together with allelic variation in the dhurrin cluster, could change allocation and the precursor pool available for lipid synthesis.

Thus, under CTL, tyrosine catabolism is enriched across independent loci, while under LIN, the dhurrin route relies on cluster members within a single block. Both implicate the same branch point, so these data motivate follow-up on tyrosine allocation between catabolism and dhurrin rather than proving a difference.

